% source: RKF/TERMINOLOGY_AND_FILING_ALIGNMENT.md
\hypertarget{terminology-and-filing-alignment}{%
\chapter{Terminology and Filing Alignment}\label{terminology-and-filing-alignment}}

\hypertarget{controlled-public-hierarchy}{%
\section{Controlled public hierarchy}\label{controlled-public-hierarchy}}

\begin{verbatim}
Recognition Kernel Framework
    public umbrella architecture

Recognition–Null Kernel Engine (RNKE)
    operational state / recognition / evolution / kernel evaluator

Recognition Residue Calculus
    parent recognition, residue, compensation, and memory logic

Recognition-Seam Calculus (RSC)
    mathematical calculus of cut, seam, recognition, memory, and open residue

EMK
    geometric and tensor recognition representation

UGD
    phase–scale–seam computational and ledger representation

RMG
    rotational, flux, motion, and dynamic representation

ECL
    executable certificate and classification runtime

Mathematical proof-transport cartridge
    current reviewer implementation imported from RH-Framework
\end{verbatim}

\hypertarget{rnke-domain-cartridge}{%
\section{RNKE domain cartridge}\label{rnke-domain-cartridge}}

A domain cartridge is represented by

\[
(S_d,R_d,E_d,\tau_d,K_d),
\]

with:

\begin{itemize}
\tightlist
\item
  \texttt{S\_d}: domain or evidence state;
\item
  \texttt{R\_d}: recognition operator, predicate, invariant, or admissibility rule;
\item
  \texttt{E\_d}: evolution, transformation, theorem step, measurement, or update;
\item
  \texttt{tau\_d}: clock, index, dependency order, or execution sequence;
\item
  \texttt{K\_d}: kernel evaluator or recognition-gap computation.
\end{itemize}

\hypertarget{mathematical-proof-transport-mapping}{%
\section{Mathematical proof-transport mapping}\label{mathematical-proof-transport-mapping}}

\begin{longtable}[]{@{}
  >{\raggedright\arraybackslash}p{(\columnwidth - 2\tabcolsep) * \real{0.5000}}
  >{\raggedright\arraybackslash}p{(\columnwidth - 2\tabcolsep) * \real{0.5000}}@{}}
\toprule\noalign{}
\begin{minipage}[b]{\linewidth}\raggedright
RNKE term
\end{minipage} & \begin{minipage}[b]{\linewidth}\raggedright
Current mathematical cartridge
\end{minipage} \\
\midrule\noalign{}
\endhead
\bottomrule\noalign{}
\endlastfoot
Domain state & theorem source, implementation, specification, and result artifacts \\
Recognition target & declared theorem obligation and dependency policy \\
Evolution & one derivation or promotion step \\
Lawful transport & proved prior theorem, exact identity, declared bound, or admitted transformation \\
Open residue & failed identity, source mismatch, missing tail, invalid rearrangement, undeclared dependency, or unclosed domain obligation \\
Kernel evaluator & deterministic verifier and negative-control suite \\
Classification & PASS, INCONCLUSIVE, FAIL, or STALE \\
Technical action & accept, repair, downgrade, rerun, or reject \\
Certificate & canonical JSON with source pins, obligations, results, and hashes \\
\end{longtable}

\hypertarget{preferred-public-wording}{%
\section{Preferred public wording}\label{preferred-public-wording}}

Use precise technical expressions such as:

\begin{verbatim}
undeclared dependency
unproved external theorem dependency
provenance mismatch
open residue
lawful transport
recognition-preserving transition
bounded residue
actionable residue
computational audit
proof-transport certificate
\end{verbatim}

Avoid loaded, adversarial, or promotional rhetoric. Public claims should describe the exact dependency, residue, failed obligation, or certificate boundary instead of assigning moral or dramatic labels to mathematical objects.

\hypertarget{historical-identifiers}{%
\section{Historical identifiers}\label{historical-identifiers}}

Certified files retain historical paths, namespaces, schemas, and package IDs. In particular, \texttt{rh\_framework} remains inside the imported cartridge because its source blob is pinned by the archived certificate. Historical identity is provenance, not the current public umbrella name.

\hypertarget{filing-alignment-notice}{%
\section{Filing-alignment notice}\label{filing-alignment-notice}}

This terminology file records technical consistency with the inventor's existing recognition-kernel filing and research records. It does not reproduce confidential filing materials, define legal claim scope, or provide an opinion on patent validity, priority, infringement, or enforceability.

\begin{flushright}\small\texttt{RKF/TERMINOLOGY\_AND\_FILING\_ALIGNMENT.md}\end{flushright}
